% chapters/02-literature-review.tex
\chapter{การทบทวนวรรณกรรม}
\label{ch:litreview}

บทนี้ทบทวนวรรณกรรมที่เกี่ยวข้องใน 7 มิติหลัก ได้แก่ (1) นิยามและการพยากรณ์
ความนิยมของเนื้อหาบนสื่อออนไลน์ (2) คุณลักษณะของวิดีโอที่ส่งผลต่อความไวรัล
(3) สถาปัตยกรรมทรานส์ฟอร์เมอร์ (4) แบบจำลองภาษาไทยที่ใช้ในงานวิจัย
(5) การจัดการความไม่สมดุล (6) การปรับเทียบความน่าจะเป็น และ
(7) ความสามารถในการอธิบายของแบบจำลอง

\section{การพยากรณ์ความนิยมของเนื้อหาบนสื่อออนไลน์}
\label{sec:lit-virality}

\subsection{นิยามและการวัดความไวรัล}

แนวคิดเรื่อง ``ความไวรัล'' ของเนื้อหาออนไลน์เริ่มต้นจาก \textcite{gladwell2000tipping} ซึ่งใช้คำนี้ในบริบทของจุดเปลี่ยน (tipping point)
ของการแพร่กระจายทางสังคม \textcite{berger2012makes} วิเคราะห์
บทความข่าวกว่า 7{,}000 ชิ้นจาก The New York Times พบว่าเนื้อหาที่กระตุ้น
\emph{อารมณ์ตื่นเต้น} (high-arousal emotion) ไม่ว่าจะเชิงบวกหรือเชิงลบมีโอกาส
ถูกแชร์สูงกว่า งานนี้ได้สถาปนาความเชื่อมโยงระหว่างคุณลักษณะของข้อความและ
การแพร่กระจาย ในงานพยากรณ์ \textcite{cheng2014can} ระบุว่าการพยากรณ์
ความสำเร็จของเนื้อหา \emph{ก่อน} เผยแพร่ทำได้ยากกว่าการพยากรณ์ \emph{หลัง}
เผยแพร่อย่างมีนัยสำคัญ เพราะสัญญาณการแพร่กระจายในระยะแรกหายไป

นิยามเชิงปริมาณของความไวรัลที่ใช้กันอย่างแพร่หลายแบ่งได้เป็น 2 แนวคือ
(ก) \emph{เกณฑ์เชิงสัมบูรณ์} เช่น \textcite{trzcinski2017popularity}
ตัดที่ 1 ล้านยอดเข้าชม และ (ข) \emph{เกณฑ์เชิงสัมพัทธ์} เช่น \textcite{khosla2014popularity} ใช้ z-score ภายในผู้สร้างเนื้อหา การศึกษานี้เลือก
แนวทาง (ข) เพราะเกณฑ์สัมบูรณ์ทำให้ช่องขนาดใหญ่กลายเป็นภาคชั้นบวกทั้งหมด ซึ่ง
ทำลายความสามารถในการแยกแยะของแบบจำลอง

\subsection{งานพยากรณ์ความนิยมบน YouTube และแพลตฟอร์มอื่น}

\textbf{\textcite{pinto2013using}} เสนอแบบจำลอง Multivariate Linear
และ Multi-Scale Series สำหรับพยากรณ์ยอดเข้าชม YouTube โดยอาศัยพฤติกรรมการ
รับชมในช่วง 30 วันแรก งานนี้บุกเบิกแนวคิด \emph{early-stage prediction} ที่
ยังคงเป็นพื้นฐานสำคัญแต่ \emph{ต้องการ} ข้อมูลหลังเผยแพร่

\textbf{\textcite{hoiles2017engagement}} ศึกษาว่าฟีเจอร์ใดบนวิดีโอ
YouTube ส่งผลต่อการมีส่วนร่วม โดยใช้ Random Forest บน meta-data รวมกับฟีเจอร์
ของช่อง สรุปว่าจำนวนผู้ติดตามและความถี่ในการอัปโหลดเป็นปัจจัยที่สำคัญที่สุด
สอดคล้องกับการค้นพบของงานปัจจุบันที่ฟีเจอร์ \texttt{channel\_subscribers}
และ \texttt{channel\_upload\_frequency} ติดอันดับสูงใน SHAP

\textbf{\textcite{trzcinski2017popularity}} ใช้คุณลักษณะภาพและ
ข้อมูลโครงสร้างของวิดีโอ Facebook บน Support Vector Regression รายงาน
$R^2 = 0.65$ ใน 24 ชั่วโมงแรก งานนี้ไม่ได้ใช้ข้อมูลข้อความเชิงลึก

\textbf{\textcite{lee2010approach}} ใช้ Naïve Bayes บนคำสำคัญของชื่อ
วิดีโอ YouTube พบว่าคำที่กระตุ้นอารมณ์มีน้ำหนักสูงในการพยากรณ์ความนิยม

\textbf{\textcite{cheng2014can}} วิเคราะห์การ ``cascade'' ของรูปภาพ
บน Facebook เน้นว่าการพยากรณ์ในขั้นแรก ๆ มีข้อจำกัดเชิงข้อมูลและความสุ่ม
สูง รูปแบบเดียวกันถูกสังเกตในข้อมูลภาษาอื่น ๆ รวมถึงเอเชียตะวันออกเฉียงใต้
เมื่อนำเทคนิคเดียวกันไปประยุกต์

ในบริบทภาษาไทย งานที่ใกล้ที่สุดคือชุดวิเคราะห์ engagement บนสื่อสังคมไทย
ที่อาศัย hand-crafted features ผสมกับ classical ML เช่น Random Forest และ
Gradient Boosting แต่ยังไม่มีงานที่ใช้แบบจำลองทรานส์ฟอร์เมอร์ภาษาไทย
หลายตัวเปรียบเทียบกันบนปัญหา virality classification พร้อมการทดสอบทาง
สถิติที่เข้มงวด งานวิจัยปัจจุบันจึงเข้าไปอุดช่องว่างเชิงระเบียบวิธีวิจัย
ดังกล่าว และเป็นการศึกษาแรก (ตามที่คณะผู้วิจัยทราบ) ที่รายงาน McNemar
pairwise + Cochran's Q ระหว่าง Thai encoders บนชุดทดสอบเดียวกัน

\section{คุณลักษณะของวิดีโอที่ส่งผลต่อความไวรัล}
\label{sec:lit-features}

\subsection{ฟีเจอร์เชิงข้อความ (Textual Features)}

ชื่อวิดีโอเป็นจุดที่ผู้ชมพบเห็นเป็นอันดับแรก (first impression)
\textcite{potthast2018clickbait} เสนอเครื่องมือ
ตรวจจับคลิกเบตที่ใช้สถาปัตยกรรม LSTM และ Transformer พบว่าฟีเจอร์ความยาว
ชื่อเรื่อง สัดส่วนอักษรพิมพ์ใหญ่ และการมีเครื่องหมายอัศเจรีย์เป็นตัวทำนายที่
สำคัญ ในงานปัจจุบันก็ยืนยันการค้นพบนี้บนข้อมูลภาษาไทย: ฟีเจอร์ \texttt{t\_emoji\_n},
\texttt{t\_question\_n}, \texttt{t\_caps\_ratio} มีอิทธิพลบน LightGBM-best
และปรากฏใน top-50 รายการ SHAP (ดูภาคผนวก \ref{app:tables})

สำหรับภาษาไทย \textcite{phatthiyaphaibun2016thainlp}
จัดทำ PyThaiNLP toolkit ที่รวมการตัดคำด้วย \texttt{newmm} ซึ่งเป็นพื้นฐาน
ของการสกัดฟีเจอร์ลำดับคำในงานนี้

\subsection{ฟีเจอร์เชิงอารมณ์ (Sentiment Features)}

\textcite{suriyawongkul2019wisesight} เสนอชุดข้อมูล
Wisesight-Sentiment ที่จัดทำขึ้นเพื่อการวิเคราะห์ความรู้สึกภาษาไทย แบ่งเป็น
4 ระดับ (positive / neutral / negative / question) งานวิจัยจำนวนมากใช้ชุดข้อมูล
นี้ฝึกตัวจำแนกความรู้สึกระดับประโยค \cite{lowphansirikul2021wangchanberta}
ในงานปัจจุบัน เราใช้ \texttt{wangchanberta-base-att-spm-uncased} ที่ปรับจูนบน
Wisesight-Sentiment แล้ว เพื่อสกัดฟีเจอร์อารมณ์ของชื่อวิดีโอ จากนั้นนำเข้า
เป็น 6 มิติของแบบจำลองคลาสสิก

\subsection{ฟีเจอร์เชิงโครงสร้าง (Channel and Temporal Features)}

ขนาดของช่องและความถี่ในการอัปโหลดเป็นตัวแปรที่ \textcite{hoiles2017engagement} และ \textcite{borghol2012untold}
ยืนยันว่ามีผลสำคัญ Borghol et~al.\ ค้นพบว่าหลักการ ``\textit{The Rich Get
Richer}'' ทำงานบน YouTube ในแง่ที่ว่าวิดีโอจากช่องขนาดใหญ่ที่มีฐานผู้ติดตาม
อยู่แล้วจะได้รับ yt-discovery bonus ส่งผลให้มีโอกาสไวรัลสูงกว่าเชิงสถิติ

ฟีเจอร์เชิงเวลา (เช่น เวลาที่เผยแพร่ในรอบสัปดาห์ ฤดูกาล วันเทศกาล) ก็มีผลตาม
\textcite{wu2016unfolding} ที่วิเคราะห์การมีส่วนร่วมตามเวลาเผยแพร่บน
Twitter

\section{แบบจำลองทรานส์ฟอร์เมอร์}
\label{sec:lit-transformer}

\subsection{สถาปัตยกรรม Transformer}

\textbf{\textcite{vaswani2017attention}} เสนอแบบจำลอง Transformer
ในปี ค.ศ.\ 2017 ซึ่งทดแทน RNN และ LSTM ในการประมวลผลลำดับด้วย
\emph{Self-Attention Mechanism} หัวใจของสถาปัตยกรรมคือฟังก์ชัน Scaled
Dot-Product Attention:
\begin{equation}
\mathrm{Attention}(Q, K, V) = \mathrm{softmax}\!\left(\frac{Q K^{\top}}{\sqrt{d_k}}\right) V
\label{eq:attention}
\end{equation}
โดย $Q, K, V$ คือ Query, Key, Value matrices และ $d_k$ คือมิติของ key
สถาปัตยกรรมเต็มของ Transformer ใช้ \emph{Multi-Head Attention} ที่รันหลาย
attention heads ขนานกัน และใช้ \emph{Positional Encoding} เพื่อแทนตำแหน่ง
ของแต่ละ token ในลำดับ จุดแข็งหลักของสถาปัตยกรรมนี้คือการเชื่อมตำแหน่งใด ๆ
สองตำแหน่งในลำดับด้วยเส้นทางคำนวณเดียว (path length $O(1)$) ซึ่งหลีกเลี่ยง
ปัญหา \emph{การเลือนหายของเกรเดียนต์} (vanishing gradient) ที่ RNN/LSTM
ประสบเมื่อต้องเรียนรู้บริบทระยะยาว แลกมาด้วยความซับซ้อนของหน่วยความจำและ
เวลาที่ $O(n^2)$ ตามความยาวลำดับ

\subsection{BERT และ RoBERTa}

\textbf{\textcite{devlin2019bert}} เสนอ BERT (Bidirectional
Encoder Representations from Transformers) ในปี ค.ศ.\ 2019 BERT ใช้
สถาปัตยกรรม Transformer-encoder อย่างเดียว ฝึกล่วงหน้าด้วยภาระงาน 2 แบบคือ
\emph{Masked Language Modelling} (MLM) ซึ่งสุ่มซ่อน token แล้วฝึกแบบจำลอง
ให้พยากรณ์ token นั้น และ \emph{Next Sentence Prediction} (NSP) เพื่อเรียนรู้
ความสัมพันธ์ระหว่างประโยค จุดแข็งของ BERT คือ \emph{สองทิศทาง}
(bidirectional) ในการเข้ารหัสบริบท ตรงกันข้ามกับ GPT ที่เป็น unidirectional
left-to-right

\textbf{\textcite{liu2019roberta}} พัฒนา RoBERTa (Robustly Optimised
BERT Pretraining Approach) ในปี ค.ศ.\ 2019 ซึ่งปรับขั้นตอนการฝึกล่วงหน้าของ
BERT ใน 4 ประเด็นหลัก คือ (1) ฝึกบนคลังข้อมูลที่ใหญ่กว่า 10 เท่า (2) ฝึกด้วย
batch ที่ใหญ่กว่า (3) ใช้ dynamic masking และ (4) ยกเลิก NSP ผลคือ RoBERTa
ทำงานได้ดีกว่า BERT ในเกือบทุก benchmark แบบจำลองทั้ง 3 ตัวที่เปรียบเทียบใน
งานนี้ใช้สถาปัตยกรรม RoBERTa

\subsection{แบบจำลองทรานส์ฟอร์เมอร์ภาษาไทย}
\label{sec:lit-thai-tx}

\textbf{WangchanBERTa.} \cite{lowphansirikul2021wangchanberta} เป็น RoBERTa
ที่ฝึกล่วงหน้าโดย AIResearch.in.th (สังกัด VISTEC) บนคลังข้อมูลภาษาไทยขนาด
78.5 GB ครอบคลุม Twitter, Wikipedia, ข่าวออนไลน์ และกฎหมาย มีพารามิเตอร์
105M ใช้ตัวตัดคำ SentencePiece ขนาดคำศัพท์ 25{,}000 token แบบจำลองนี้เป็นที่
รู้จักในชื่อ \texttt{airesearch/wangchanberta-base-att-spm-uncased} และ
ได้กลายเป็นพื้นฐาน (baseline) มาตรฐานของงาน NLP ภาษาไทย

\textbf{PhayaThaiBERT.} \cite{phayathaibert2023} เป็นแบบจำลองภาษาไทยที่ใหญ่กว่า
พัฒนาโดยทีม CLICKNEXT ฝึกบนคลังข้อมูล 312 GB ขนาดพารามิเตอร์ 278M ใช้ตัวตัด
คำ SentencePiece เช่นกันแต่มีขนาดคำศัพท์ 250{,}000 token ซึ่งครอบคลุมคำเฉพาะ
ทางและคำใหม่ในภาษาไทยปัจจุบันได้ดีกว่า งานวิจัยภาษาไทยรายงานว่า PhayaThaiBERT
ทำงานได้ดีกว่า WangchanBERTa บนงาน sentiment, NER และ Thai understanding
benchmarks

\textbf{XLM-RoBERTa-large.} \cite{conneau2020xlmr} เป็น RoBERTa ที่ฝึกแบบ
มัลติลิงกัวลบนคลังข้อมูล 100 ภาษา รวม 2.5 TB จาก CommonCrawl มีพารามิเตอร์
560M ใช้ SentencePiece ขนาดคำศัพท์ 250{,}000 token ครอบคลุมภาษาไทย ภาษา
อังกฤษ และภาษาในกลุ่มประเทศอาเซียน สถาปัตยกรรมรองรับการสลับภาษาในเอกสาร
เดียวกัน ในทฤษฎี XLM-R-large ควรได้เปรียบบน YouTube ภาษาไทยที่มักปะปนกับ
อังกฤษ แต่ผลการทดลองของงานนี้กลับขัดกับสมมติฐาน

\subsection{การปรับจูน (Fine-tuning) แบบเต็มเทียบกับ Parameter-Efficient}

\textbf{\textcite{howard2018ulmfit}} แสดงให้เห็นว่าการปรับจูนแบบ
จำลองภาษาที่ผ่านการฝึกล่วงหน้าให้ดีกว่าการฝึกใหม่จากศูนย์ โดยเฉพาะเมื่อข้อมูล
มีจำกัด การปรับจูนแบบเต็ม (full fine-tuning) ปรับน้ำหนักทั้งหมดของแบบจำลอง
ซึ่งทำงานได้ดีที่สุดเมื่อทรัพยากรเพียงพอ ในงานนี้เราใช้ full FT เพราะข้อมูล
และ GPU มีเพียงพอ

ทางเลือกอื่น ๆ คือ \emph{Parameter-Efficient Fine-Tuning} (PEFT) เช่น
\textcite{hu2022lora} และ \textcite{dettmers2023qlora} ซึ่งฝึกเฉพาะ
adapter ขนาดเล็ก ลดความต้องการ VRAM ลงได้มาก เหมาะกับแบบจำลองขนาดใหญ่ระดับ
หลายพันล้านพารามิเตอร์ รหัสฐานของงานนี้รองรับ QLoRA สำหรับ decoder LLM ผ่าน
\texttt{src/models/transformer\_finetune.py:\_build\_trainer\_qlora()} แต่ไม่ได้
ใช้ในการเปรียบเทียบหลัก

\section{แบบจำลองพื้นฐานเชิงคลาสสิกและการเรียนรู้เชิงผสม}
\label{sec:lit-classical}

\subsection{Logistic Regression, LightGBM, XGBoost}

\textbf{Logistic Regression} เป็น \emph{แบบจำลองพื้นฐาน} (baseline) ที่นิยมใช้
อย่างแพร่หลายสำหรับการจำแนกประเภทไบนารี \textbf{XGBoost} \cite{chen2016xgboost}
และ \textbf{LightGBM} \cite{ke2017lightgbm} เป็น Gradient Boosting Decision
Tree รุ่นที่ออกแบบมาเพื่อความเร็วในการฝึกและจัดการกับข้อมูลขนาดใหญ่ได้ดี
LightGBM ใช้ leaf-wise tree growth และ histogram-based split finding ซึ่ง
เร็วกว่า XGBoost ในกรณีส่วนใหญ่ ในงาน benchmark Kaggle, LightGBM และ XGBoost
มักจะอยู่ใน top-3 ของแบบจำลองที่ใช้บนข้อมูล tabular

\subsection{TF-IDF กับการจำแนกข้อความ}

\textbf{Sparck \textcite{sparckjones1972tfidf}} เสนอ Inverse Document
Frequency (IDF) ในปี ค.ศ.\ 1972 ซึ่งกลายเป็นพื้นฐานของ TF-IDF การใช้
character-level n-gram TF-IDF กับภาษาไทย ที่ไม่มีการเว้นวรรค แสดงให้เห็นว่า
ทำงานได้ดีกว่า word-level เนื่องจากปัญหา OOV (out-of-vocabulary) และความ
ไม่แน่นอนของขอบเขตคำ

\subsection{แบบจำลองไฮบริด (Hybrid Models)}

แบบจำลองไฮบริดเป็นวิธีการรวมข้อมูลหลายโมดาลิตี (text + structured) ที่ได้รับ
ความนิยมในงาน Kaggle เช่น KDD Cup 2014 และในอุตสาหกรรม \textcite{cheng2016widedeep} เสนอ Wide \& Deep โดย Google ที่รวม linear features
กับ embeddings เข้าด้วยกัน งานปัจจุบันใช้แนวคิดเดียวกัน แต่แทน wide ด้วย
LightGBM head และแทน deep ด้วย MLP บน title embeddings จาก SBERT

\subsection{การรวมแบบจำลองแบบเรียงซ้อน (Stacking Ensemble)}

\textbf{Wolpert \cite{wolpert1992stacked}} เสนอแนวคิด stacking generalisation
ในปี ค.ศ.\ 1992 ซึ่งใช้ \emph{แบบจำลองระดับเมตา} (meta-learner) เรียนรู้วิธี
รวมผลพยากรณ์ของแบบจำลองฐาน (base models) หลายตัว \textcite{breiman1996stacked}
ขยายและพิสูจน์ว่า stacking ลด generalisation error ลงตามทฤษฎี การประยุกต์ใช้
stacking ที่ประสบความสำเร็จมีมากในการประกวด Kaggle เช่น Netflix Prize
\textcite{bell2007lessons} รายงานว่าการรวมแบบจำลองหลายตัวช่วยเพิ่ม
ประสิทธิภาพอย่างมีนัยสำคัญ ในงานปัจจุบันเราใช้ Logistic Regression เป็น
meta-learner เพราะตีความได้ง่ายและไม่ overfit ง่ายเมื่อแบบจำลองฐานเป็นความ
น่าจะเป็นที่ปรับเทียบมาแล้ว

\section{การจัดการความไม่สมดุล}
\label{sec:lit-imbalance}

ในปัญหาปัจจุบันสัดส่วนภาคชั้นบวกอยู่ที่ 10.16\% ซึ่งถือว่ามีความไม่สมดุล
ปานกลาง (moderate imbalance) วิธีแก้ไขที่ใช้กันทั่วไปแบ่งเป็น 3 กลุ่ม

\subsection{การปรับสมดุลของข้อมูล (Resampling)}

\textbf{Undersampling} ลดจำนวนตัวอย่างของชั้นที่มีข้อมูลมาก (majority class —
ในที่นี้คือชั้นลบ) ลงให้สมดุลกับชั้นที่มีข้อมูลน้อย (minority class — ชั้นบวก)
ข้อจำกัดของวิธีนี้คือสูญเสียข้อมูลส่วนหนึ่ง \textbf{Oversampling} เช่น \textcite{chawla2002smote} สร้างตัวอย่างของชั้นที่มีข้อมูลน้อยเพิ่มโดยสังเคราะห์
ระหว่างเพื่อนบ้าน $k$-NN ข้อจำกัดคือใน feature space ที่มีมิติสูงและกระจัด
กระจาย SMOTE อาจสร้างข้อมูลที่ไม่สมจริง \textcite{he2009learning}
สรุปว่าสำหรับ NLP บนข้อความสั้น undersampling มักเป็นทางเลือกที่ปลอดภัยกว่า

\subsection{การถ่วงน้ำหนัก (Class Weighting) และ Focal Loss}

วิธีที่ตรงไปตรงมาคือการถ่วงน้ำหนักให้กับภาคชั้นที่มีข้อมูลน้อยใน cross-entropy
loss \textbf{Focal Loss} \cite{lin2017focal} ถูกเสนอเพื่อให้แบบจำลองเน้น
ตัวอย่างที่จำแนกได้ยาก โดยใส่แฟกเตอร์ $(1-p_t)^\gamma$ เข้าไปลดน้ำหนักของ
ตัวอย่างที่จำแนกได้ง่าย ในงานวิจัยปัจจุบันการทดลองเบื้องต้นบน WangchanBERTa
พบว่า Focal Loss ($\gamma=2$) ไม่ได้ให้ผลดีกว่า cross-entropy แบบธรรมดาที่
รวมกับ undersampling

\subsection{การปรับเปลี่ยนค่าเกณฑ์ตัดสิน (Decision Threshold)}

นอกจากปรับ training-side ยังสามารถปรับ inference-side ด้วยการเลื่อน
ค่าเกณฑ์การจำแนกออกจาก 0.5 ไปยังจุดที่เหมาะสมที่สุดสำหรับ F1 ของชั้นบวก
(\textcite{he2008imbalanced}) ในงานวิจัยปัจจุบันเราใช้แนวทางนี้:
ค่าเกณฑ์ถูกเลือกจากชุดตรวจสอบด้วยการสแกนอย่างละเอียด แล้วใช้เป็น
ค่าเกณฑ์ในการอนุมานสุดท้าย

\section{การปรับเทียบความน่าจะเป็น}
\label{sec:lit-calibration}

\subsection{Platt Scaling}

\textbf{Platt \cite{platt1999platt}} เสนอ Platt scaling ในปี ค.ศ.\ 1999
เพื่อปรับเทียบคะแนน decision function ของ SVM ให้กลายเป็นความน่าจะเป็น
หลักการคือฝึก Logistic Regression ตัวที่สองบน score ที่ได้
\begin{equation}
P(y=1\mid s) = \frac{1}{1 + \exp(As + B)}
\label{eq:platt}
\end{equation}
โดย $s$ คือคะแนนพยากรณ์ ตัวแปร $A, B$ ปรับด้วย MLE บน validation fold
Platt scaling เหมาะกับการบิดของ sigmoid ที่อาจเกิดจาก SVM และโครงข่าย
ประสาทเทียม

\subsection{Isotonic Regression}

\textbf{\textcite{zadrozny2002isotonic}} เสนอ isotonic regression
ในงาน calibration ซึ่งเป็น non-parametric piecewise constant monotone
mapping ข้อดีคือยืดหยุ่นกว่า Platt และไม่จำกัดให้อยู่ในรูป sigmoid ข้อเสียคือ
ต้องการข้อมูล validation เยอะกว่าและ overfit ง่ายเมื่อข้อมูลน้อย

\subsection{การวัดคุณภาพการปรับเทียบ}

Niculescu-\textcite{niculescu2005predicting} เปรียบเทียบ
calibration ของแบบจำลองต่าง ๆ พบว่า boosted trees มักให้ความน่าจะเป็นที่
\emph{เอียง} ไปทาง 0.5 ส่วน RF และ NB เอียงไปทาง 0 หรือ 1 \emph{Expected
Calibration Error} (ECE) \cite{naeini2015obtaining} เป็นตัววัดมาตรฐาน:
\begin{equation}
\mathrm{ECE} = \sum_{m=1}^{M} \frac{|B_m|}{n}\,\bigl|\mathrm{acc}(B_m) - \mathrm{conf}(B_m)\bigr|
\label{eq:ece}
\end{equation}
โดยแบ่งข้อมูลเป็น $M$ bins ตาม confidence แล้ววัดส่วนต่างของความแม่นยำกับ
ความมั่นใจในแต่ละ bin งานวิจัยปัจจุบันรายงาน ECE ของแบบจำลองที่ดีที่สุด
หลังการปรับเทียบที่ 0.016 ซึ่งถือว่าดีมาก

\section{การทดสอบความแตกต่างทางสถิติของแบบจำลอง}
\label{sec:lit-stats}

\subsection{McNemar's Test}

\textbf{McNemar \cite{mcnemar1947}} เสนอการทดสอบ paired binary outcomes
ในปี ค.ศ.\ 1947 ในการเปรียบเทียบ 2 แบบจำลองบนชุดทดสอบเดียวกัน เราสร้างตาราง
$2\times 2$:
\begin{center}
\begin{tabular}{l|cc}
 & B ถูก & B ผิด \\\hline
A ถูก & $n_{00}$ & $n_{01}$ \\
A ผิด & $n_{10}$ & $n_{11}$
\end{tabular}
\end{center}
สถิติของ McNemar คือ
\begin{equation}
\chi^2_{\mathrm{McNemar}} = \frac{(n_{01} - n_{10})^2}{n_{01} + n_{10}}
\label{eq:mcnemar}
\end{equation}
ภายใต้ $H_0$ ที่ว่าแบบจำลอง 2 ตัวมีอัตราความผิดพลาดเท่ากัน สถิติเข้าใกล้
การแจกแจง $\chi^2_1$ \textcite{dietterich1998approximate}
แนะนำการทดสอบนี้สำหรับการเปรียบเทียบ 2 ตัวจำแนก

\subsection{Cochran's Q Test}

\textbf{Cochran \cite{cochran1950}} ขยาย McNemar ไปยังการเปรียบเทียบ $k > 2$
ตัวจำแนก สถิติคือ
\begin{equation}
Q = \frac{(k-1)\,\left[k \sum_{j=1}^{k} C_j^2 - T^2\right]}{kT - \sum_{i=1}^{n} R_i^2}
\label{eq:cochranq}
\end{equation}
โดย $C_j$ คือผลรวมความถูกของตัวจำแนก $j$ ทุกตัวอย่าง $R_i$ คือผลรวมความถูก
ของตัวอย่าง $i$ ข้ามตัวจำแนก และ $T$ คือผลรวมทั้งหมด ภายใต้ $H_0$ สถิติเข้า
ใกล้ $\chi^2_{k-1}$ งานวิจัยปัจจุบันใช้ Cochran's Q กับทั้ง 22 แบบจำลองและ
กับสามทรานส์ฟอร์เมอร์ตามเงื่อนไขแยก

\subsection{Bootstrap Confidence Intervals}

\textcite{efron1979bootstrap} เสนอ Bootstrap resampling สำหรับประมาณการ
sampling distribution โดยไม่ต้องสมมติฐาน parametric งานวิจัยปัจจุบันใช้
1{,}000-sample bootstrap แบบ stratified โดย resample บนชุดทดสอบเพื่อสร้างช่วง
ความเชื่อมั่น 95\% ของ ROC-AUC วิธี percentile

\section{ความสามารถในการอธิบายของแบบจำลอง}
\label{sec:lit-xai}

\subsection{SHAP}

\textbf{\textcite{lundberg2017shap}} เสนอ SHAP (SHapley Additive
exPlanations) ในปี ค.ศ.\ 2017 ซึ่งใช้ Shapley value จากทฤษฎีเกมเพื่อกระจาย
แรงโน้มน้าวของแต่ละฟีเจอร์ต่อพยากรณ์ของแบบจำลอง คุณสมบัติเด่นคือ
\emph{additivity} และ \emph{local accuracy} ทำให้ผลรวม SHAP ของทุกฟีเจอร์
เท่ากับส่วนต่างของพยากรณ์จากค่ากลาง การคำนวณ Shapley value โดยตรงเป็น NP-hard
แต่ SHAP มีการประมาณการที่มีประสิทธิภาพสำหรับ tree-based model (TreeSHAP)
และโครงข่ายประสาทเทียม (DeepSHAP, GradientSHAP)

\subsection{LIME}

\textbf{\textcite{ribeiro2016should}} เสนอ LIME (Local
Interpretable Model-Agnostic Explanations) ซึ่งสร้างคำอธิบายเฉพาะที่
(local) โดยฝึก surrogate model เชิงเส้นบนตัวอย่างที่ perturb รอบ ๆ จุดที่
สนใจ ในงาน NLP perturbation ทำโดย ``ลบ'' คำออกจากประโยคแล้วประเมินการ
เปลี่ยนแปลงของพยากรณ์ จุดเด่นคือ \emph{model-agnostic} (ใช้กับแบบจำลองอะไร
ก็ได้) จุดด้อยคือเสถียรภาพต่ำเมื่อตัวอย่าง perturbation ไม่มากพอ

\subsection{Attention Rollout}

\textbf{\textcite{abnar2020quantifying}} เสนอ ``attention
rollout'' ในปี ค.ศ.\ 2020 เป็นวิธีรวม attention weights ข้าม layer ของ
Transformer เป็นความสนใจรวมจาก \texttt{[CLS]} token ไปยัง input tokens
หลักการคือคูณ attention matrices ข้าม layer หลังเพิ่ม identity ตัวแทน
residual:
\begin{equation}
\tilde A^{(l)} = \tfrac{1}{2}\bigl(A^{(l)} + I\bigr),
\qquad
\mathrm{Rollout}^{(L)} = \prod_{l=1}^{L} \tilde A^{(l)}
\label{eq:attnroll}
\end{equation}
ผลลัพธ์คือ ``ความสนใจรวม'' ของแต่ละ token ที่มีต่อการพยากรณ์สุดท้าย
\textcite{chefer2021transformer} ขยายแนวคิดนี้ด้วย gradient-weighted
relevance propagation งานวิจัยปัจจุบันใช้ rollout แบบ Abnar \& Zuidema เพราะ
เป็น baseline ที่ใช้แพร่หลายและไม่ต้อง backprop

\section{งานวิจัยที่เกี่ยวข้อง}
\label{sec:lit-summary}

วรรณกรรมที่ทบทวนชี้ให้เห็นว่า (1) การพยากรณ์ความไวรัลก่อนเผยแพร่เป็นโจทย์ที่
ทำได้ยากเชิงข้อมูล (2) การใช้แบบจำลองภาษาขนาดใหญ่ที่ฝึกล่วงหน้าเป็น state-of-the-art
ปัจจุบันสำหรับ NLP ภาษาไทย (3) Stacking และการปรับเทียบความน่าจะเป็นเป็น
เทคนิคที่ผ่านการพิสูจน์แล้วในการเพิ่มประสิทธิภาพ (4) การทดสอบทางสถิติเช่น
McNemar และ Cochran's Q เป็นมาตรฐานของการเปรียบเทียบ classifiers และ (5)
SHAP/LIME/Attention rollout เป็นเครื่องมือ XAI ที่ได้รับการยอมรับ
งานวิจัยปัจจุบันสร้างจุดร่วมระหว่างทั้ง 5 มิตินี้ในบริบทภาษาไทยที่ยังไม่มี
งานก่อนหน้าทำอย่างครบถ้วน
